% questions.tex —— 2014 年·上海卷·理科 · 题目内容片段（无答案）
% 由 tools/generate_projects.py 依来源总账自 2014/2014_数学真题_上海理_试题.md 生成，勿手改（改源 MD 后重跑）


\begin{gknotice}
\textbf{来源说明}：本题面依据固定源 PDF 的两页 A3 题面逐题转录，原卷印刷页码为 915—916。第 915 面为填空题 1—14 与选择题 15—18；第 916 面为解答题 19—23。全卷不印分值，三处节头依次为「一、填空题」「二、选择题」「三、解答题」。

\textbf{解析归属}：固定源不附答案或解析；配套答案与解析均由本库按本题面独立复算。

\textbf{图形说明}：第 16 题正四棱柱、第 19 题表面展开图、第 21 题广告牌示意图均由项目内 TikZ 矢量源重绘；其余各题原页无图。
\end{gknotice}

\gkpart{一、填空题}

\begin{question}
\gkti{1}函数 $y=1-2\cos^2(2x)$ 的最小正周期是\gkblank{}。

\end{question}
\begin{question}
\gkti{2}若复数 $z=1+2\mathrm i$，其中 $\mathrm i$ 是虚数单位，则 $\left(z+\dfrac1{\overline z}\right)\cdot\overline z=$\gkblank{}。

\end{question}
\begin{question}
\gkti{3}若抛物线 $y^2=2px$ 的焦点与椭圆 $\dfrac{x^2}{9}+\dfrac{y^2}{5}=1$ 的右焦点重合，则该抛物线的准线方程为\gkblank{}。

\end{question}
\begin{question}
\gkti{4}设
\nopagebreak[4]
\[f(x)=\begin{cases}
x, & x\in(-\infty,a),\\
x^2, & x\in[a,+\infty),
\end{cases}\]
若 $f(2)=4$，则 $a$ 的取值范围为\gkblank{}。

\end{question}
\begin{question}
\gkti{5}若实数 $x$，$y$ 满足 $xy=1$，则 $x^2+2y^2$ 的最小值为\gkblank{}。

\end{question}
\begin{question}
\gkti{6}若圆锥的侧面积是底面积的 $3$ 倍，则其母线与底面夹角的大小为\gkblank{}（结果用反三角函数值表示）。

\end{question}
\begin{question}
\gkti{7}已知曲线 $C$ 的极坐标方程为 $\rho(3\cos\theta-4\sin\theta)=1$，则 $C$ 与极轴的交点到极点的距离是\gkblank{}。

\end{question}
\begin{question}
\gkti{8}设无穷等比数列 $\{a_n\}$ 的公比为 $q$，若 $a_1=\lim\limits_{n\to\infty}(a_3+a_4+\cdots+a_n)$，则 $q=$\gkblank{}。

\end{question}
\begin{question}
\gkti{9}若 $f(x)=x^{\frac23}-x^{-\frac12}$，则满足 $f(x)<0$ 的 $x$ 的取值范围是\gkblank{}。

\end{question}
\begin{question}
\gkti{10}为强化安全意识，某商场拟在未来的连续 $10$ 天中随机选择 $3$ 天进行紧急疏散演练，则选择的 $3$ 天恰好为连续 $3$ 天的概率是\gkblank{}（结果用最简分数表示）。

\end{question}
\begin{question}
\gkti{11}已知互异的复数 $a$，$b$ 满足 $ab\neq0$，集合 $\{a,b\}=\{a^2,b^2\}$，则 $a+b=$\gkblank{}。

\end{question}
\begin{question}
\gkti{12}设常数 $a$ 使方程 $\sin x+\sqrt3\cos x=a$ 在闭区间 $[0,2\pi]$ 上恰有三个解 $x_1$，$x_2$，$x_3$，则 $x_1+x_2+x_3=$\gkblank{}。

\end{question}
\begin{question}
\gkti{13}某游戏的得分为 $1,2,3,4,5$，随机变量 $\xi$ 表示小白玩该游戏的得分。若 $E(\xi)=4.2$，则小白得 $5$ 分的概率至少为\gkblank{}。

\end{question}
\begin{question}
\gkti{14}已知曲线 $C:x=-\sqrt{4-y^2}$，直线 $l:x=6$。若对于点 $A(m,0)$，存在 $C$ 上的点 $P$ 和 $l$ 上的 $Q$ 使得 $\overrightarrow{AP}+\overrightarrow{AQ}=\vec{0}$，则 $m$ 的取值范围为\gkblank{}。

\end{question}
\gkpart{二、选择题}

\begin{question}
\gkti{15}设 $a$，$b\in\mathbf R$，则“$a+b>4$”是“$a>2$ 且 $b>2$”的\mbox{（　　）}

（A）充分非必要条件；（B）必要非充分条件；（C）充要条件；（D）既不充分也不必要条件。

\Needspace{0.4\textheight}
\end{question}
\begin{question}
\gkti{16}如图，四个棱长为 $1$ 的正方体排成一个正四棱柱，$AB$ 是一条侧棱，$P_i$（$i=1,2,\cdots,8$）是上底面上其余的八个点，则 $\overrightarrow{AB}\cdot\overrightarrow{AP_i}$（$i=1,2,\cdots,8$）的不同值的个数为\mbox{（　　）}
\nopagebreak[4]
\begin{center}% ============================================================
% f2014_shanghai_science_q16_prism.tex
% 2014 上海卷（理科）第 16 题：四个棱长为 1 的正方体排成的正四棱柱
% 来源 PDF：普通高考/2014/2014上海理.pdf
% 原页：第 1 页（A3），印面 915，第 16 题题干右侧
% 底面 2×2、高 1；AB 为一条侧棱（A 在下、B 在上），
% 上底面除 B 外的八个格点为 P1—P8：
%   P1、P2 在左边（自 B 起），P3、P6 在前边，P4 为面心，
%   P5、P8 在后边，P7 在右边。
% 可见棱与上底面网格线为实线，被遮挡的棱与内部竖棱为虚线。
% 平行投影：前边方向取 1 单位 =（1.025,-0.15），竖直方向取 1 单位 =（0,-0.95）。
% ============================================================
\begin{tikzpicture}[>=Stealth, line width=0.65pt, font=\small, scale=1.1,
  line cap=round, line join=round]
  \coordinate (B) at (0,2.30);
  % 上底面九个格点
  \coordinate (P1) at ($(B)+(0.55,0.55)$);
  \coordinate (P2) at ($(B)+(1.10,1.10)$);
  \coordinate (P3) at ($(B)+(2.05,-0.30)$);
  \coordinate (P4) at ($(B)+(2.60,0.25)$);
  \coordinate (P5) at ($(B)+(3.15,0.80)$);
  \coordinate (P6) at ($(B)+(4.10,-0.60)$);
  \coordinate (P7) at ($(B)+(4.65,-0.05)$);
  \coordinate (P8) at ($(B)+(5.20,0.50)$);
  % 下底面九个格点（竖直下移 0.95）
  \coordinate (A)   at ($(B)+(0,-0.95)$);
  \coordinate (P1b) at ($(P1)+(0,-0.95)$);
  \coordinate (P2b) at ($(P2)+(0,-0.95)$);
  \coordinate (P3b) at ($(P3)+(0,-0.95)$);
  \coordinate (P4b) at ($(P4)+(0,-0.95)$);
  \coordinate (P5b) at ($(P5)+(0,-0.95)$);
  \coordinate (P6b) at ($(P6)+(0,-0.95)$);
  \coordinate (P7b) at ($(P7)+(0,-0.95)$);
  \coordinate (P8b) at ($(P8)+(0,-0.95)$);

  % 可见的外轮廓、可见棱与前、右底边
  \draw (A) -- (B) -- (P2) -- (P8) -- (P8b) -- (P6b) -- (A);
  \draw (B) -- (P6) -- (P8);
  \draw (P6) -- (P6b);
  \draw (P3) -- (P3b);
  \draw (P7) -- (P7b);

  % 上底面 2×2 网格线（过面心 P4）
  \draw (P3) -- (P5);
  \draw (P1) -- (P7);

  % 被遮挡的棱：后左竖棱、底面后边与左边、三条背面/内部竖棱
  \draw[dashed] (P2) -- (P2b);
  \draw[dashed] (P2b) -- (P8b);
  \draw[dashed] (P2b) -- (A);
  \draw[dashed] (P1) -- (P1b);
  \draw[dashed] (P4) -- (P4b);
  \draw[dashed] (P5) -- (P5b);

  % 顶点标注
  \fill (A) circle (0.85pt) node[below left=-1pt] {$A$};
  \fill (B) circle (0.85pt) node[left=-1pt] {$B$};
  \fill (P1) circle (0.85pt) node[left=-2pt] {$P_1$};
  \fill (P2) circle (0.85pt) node[above left=-1pt] {$P_2$};
  \fill (P3) circle (0.85pt) node[below left=-2pt] {$P_3$};
  \fill (P4) circle (0.85pt) node[above left=-2pt] {$P_4$};
  \fill (P5) circle (0.85pt) node[above=-1pt] {$P_5$};
  \fill (P6) circle (0.85pt) node[below left=-2pt] {$P_6$};
  \fill (P7) circle (0.85pt) node[right=-1pt] {$P_7$};
  \fill (P8) circle (0.85pt) node[above right=-1pt] {$P_8$};
\end{tikzpicture}
\end{center}
\nopagebreak[4]
（A）$1$；（B）$2$；（C）$4$；（D）$8$。

\end{question}
\begin{question}
\gkti{17}已知 $P_1(a_1,b_1)$ 与 $P_2(a_2,b_2)$ 是直线 $y=kx+1$（$k$ 为常数）上两个不同的点，则关于 $x$ 和 $y$ 的方程组
\nopagebreak[4]
\[\begin{cases}
a_1x+b_1y=1,\\
a_2x+b_2y=1
\end{cases}\]
的解的情况是\mbox{（　　）}

（A）无论 $k$，$P_1$，$P_2$ 如何，总是无解；（B）无论 $k$，$P_1$，$P_2$ 如何，总有唯一解；（C）存在 $k$，$P_1$，$P_2$，使之恰有两解；（D）存在 $k$，$P_1$，$P_2$，使之有无穷多解。

\end{question}
\begin{question}
\gkti{18}设
\nopagebreak[4]
\[f(x)=\begin{cases}
(x-a)^2, & x\leqslant0,\\
x+\dfrac1x+a, & x>0,
\end{cases}\]
若 $f(0)$ 是 $f(x)$ 的最小值，则 $a$ 的取值范围为\mbox{（　　）}

（A）$[-1,2]$；（B）$[-1,0]$；（C）$[1,2]$；（D）$[0,2]$。

\end{question}
\gkpart{三、解答题}

\Needspace{0.4\textheight}
\begin{question}
\gkti{19}底面边长为 $2$ 的正三棱锥 $P\text{-}ABC$，其表面展开图是三角形 $P_1P_2P_3$，如图，求 $\triangle P_1P_2P_3$ 的各边长及此三棱锥的体积 $V$。
\nopagebreak[4]
\begin{center}% ============================================================
% f2014_shanghai_science_q19_net.tex
% 2014 上海卷（理科）第 19 题：正三棱锥 P-ABC 的表面展开图
% 来源 PDF：普通高考/2014/2014上海理.pdf
% 原页：第 1 页（A3），印面 915，第 19 题题干下方（右栏下部）
% 展开图为三角形 P1P2P3：A、C 分别为 P1P3、P3P2 的中点，
% B 为 P1P2 的中点；△ABC 为底面（实线），大三角形三边为
% 展开图外边界（虚线）。
% ============================================================
\begin{tikzpicture}[>=Stealth, line width=0.65pt, font=\small, scale=0.78,
  line cap=round, line join=round]
  \coordinate (P1) at (0,0);
  \coordinate (P2) at (4,0);
  \coordinate (P3) at (2,3.4641);
  \coordinate (A)  at (1,1.73205);
  \coordinate (B)  at (2,0);
  \coordinate (C)  at (3,1.73205);

  % 展开图外边界（虚线）
  \draw[dashed] (P1) -- (P2) -- (P3) -- cycle;

  % 底面三角形（折痕，实线）
  \draw (A) -- (B) -- (C) -- cycle;

  % 顶点标注
  \fill (P1) circle (0.85pt) node[below left=-1pt] {$P_1$};
  \fill (P2) circle (0.85pt) node[below right=-1pt] {$P_2$};
  \fill (P3) circle (0.85pt) node[above=-1pt] {$P_3$};
  \fill (A) circle (0.85pt) node[above left=-1pt] {$A$};
  \fill (B) circle (0.85pt) node[below=-1pt] {$B$};
  \fill (C) circle (0.85pt) node[above right=-1pt] {$C$};
\end{tikzpicture}
\end{center}
\nopagebreak[4]
\end{question}
\begin{question}
\gkti{20}设常数 $a\geqslant0$，函数 $f(x)=\dfrac{2^x+a}{2^x-a}$。

（1） 若 $a=4$，求函数 $y=f(x)$ 的反函数 $y=f^{-1}(x)$；

（2） 根据 $a$ 的不同取值，讨论函数 $y=f(x)$ 的奇偶性，并说明理由。

\Needspace{0.4\textheight}
\end{question}
\begin{question}
\gkti{21}如图，某公司要在 $A$、$B$ 两地连线上的定点 $C$ 处建造广告牌 $CD$，其中 $D$ 为顶端，$AC$ 长 $35$ 米，$CB$ 长 $80$ 米。设点 $A$、$B$ 在同一水平面上，从 $A$ 和 $B$ 看 $D$ 的仰角分别为 $\alpha$ 和 $\beta$。
\nopagebreak[4]
\begin{center}% ============================================================
% f2014_shanghai_science_q21_billboard.tex
% 2014 上海卷（理科）第 21 题：广告牌 CD 的仰角示意图
% 来源 PDF：普通高考/2014/2014上海理.pdf
% 原页：第 2 页（A3），印面 916，第 21 题题干下方（左栏）
% A、C、B 在同一水平线上，AC=35 米、CB=80 米，D 为广告牌顶端；
% 图中 CD 为铅垂方向（第 (1) 问的情形），α、β 分别为从 A、B
% 看 D 的仰角。示意图按 AC:CB:CD≈35:80:28 的比例绘制。
% ============================================================
\begin{tikzpicture}[>=Stealth, line width=0.65pt, font=\small, scale=1.0,
  line cap=round, line join=round]
  \coordinate (A) at (0,0);
  \coordinate (C) at (3.50,0);
  \coordinate (B) at (11.50,0);
  \coordinate (D) at (3.50,2.80);

  % 水平线与广告牌
  \draw (A) -- (B);
  \draw (D) -- (C);
  \draw (A) -- (D);
  \draw (B) -- (D);

  % 仰角 α、β 的角标
  \draw (1.00,0) arc (0:38.66:1.00);
  \draw ([shift=(160.71:1.00)]B) arc (160.71:180:1.00);

  % 顶点标注
  \fill (A) circle (0.85pt) node[below left=-1pt] {$A$};
  \fill (C) circle (0.85pt) node[below=-1pt] {$C$};
  \fill (B) circle (0.85pt) node[below right=-1pt] {$B$};
  \fill (D) circle (0.85pt) node[above=-1pt] {$D$};
  \node at (1.50,0.45) {$\alpha$};
  \node at (10.30,0.62) {$\beta$};
\end{tikzpicture}
\end{center}
\nopagebreak[4]
（1） 设计中 $CD$ 是铅垂方向。若要求 $\alpha\geqslant2\beta$，问 $CD$ 的长至多为多少（结果精确到 $0.01$ 米）？

（2） 施工完成后，$CD$ 与铅垂方向有偏差。现在实测得 $\alpha=38.12^\circ$，$\beta=18.45^\circ$，求 $CD$ 的长（结果精确到 $0.01$ 米）。

\end{question}
\begin{question}
\gkti{22}在平面直角坐标系 $xOy$ 中，对于直线 $l:ax+by+c=0$ 和点 $P_1(x_1,y_1)$，$P_2(x_2,y_2)$，记 $\eta=(ax_1+by_1+c)(ax_2+by_2+c)$。若 $\eta<0$，则称点 $P_1$，$P_2$ 被直线 $l$ 分隔。若曲线 $C$ 与直线 $l$ 没有公共点，且曲线 $C$ 上存在点 $P_1$，$P_2$ 被直线 $l$ 分隔，则称直线 $l$ 为曲线 $C$ 的一条分隔线。

（1） 求证：点 $A(1,2)$，$B(-1,0)$ 被直线 $x+y-1=0$ 分隔；

（2） 若直线 $y=kx$ 是曲线 $x^2-4y^2=1$ 的分隔线，求实数 $k$ 的取值范围；

（3） 动点 $M$ 到点 $Q(0,2)$ 的距离与到 $y$ 轴的距离之积为 $1$，设点 $M$ 的轨迹为 $E$，求证：通过原点的直线中，有且仅有一条直线是 $E$ 的分隔线。

\end{question}
\begin{question}
\gkti{23}已知数列 $\{a_n\}$ 满足 $\dfrac13a_n\leqslant a_{n+1}\leqslant3a_n$，$n\in\mathbf N^*$，$a_1=1$。

（1） 若 $a_2=2$，$a_3=x$，$a_4=9$，求 $x$ 的取值范围；

（2） 设 $\{a_n\}$ 是公比为 $q$ 的等比数列，$S_n=a_1+a_2+\cdots+a_n$。若 $\dfrac13S_n\leqslant S_{n+1}\leqslant3S_n$，$n\in\mathbf N^*$，求 $q$ 的取值范围；

（3） 若 $a_1,a_2,\cdots,a_k$ 成等差数列，且 $a_1+a_2+\cdots+a_k=1000$，求正整数 $k$ 的最大值，以及 $k$ 取最大值时相应数列 $a_1,a_2,\cdots,a_k$ 的公差。

\end{question}